\chapter{User Guide}
\label{appendix:user-guide}

This appendix provides a step-by-step guide for using the {\name} system through its web interface.

\section{Getting Started}
\label{app:getting_started}

\subsection{Accessing the System}

After installation and startup (see Appendix~\ref{appendix:tech-specs}), access the {\name} system by navigating to \texttt{http://localhost:3000} in your web browser.

\begin{figure}[h]
\centering
\includegraphics[width=\textwidth]{ui_startup_screen.png}
\caption{Initial System Startup Screen}
\label{fig:ui_startup_screen}
\end{figure}

\section{Basic Workflow}
\label{app:basic_workflow}

\subsection{Step 1: Enter URLs}

\begin{enumerate}
\item Navigate to the Home tab (🏠 Home)
\item Enter URLs in the text area, one per line or paste multiple URLs
\item URLs will automatically be grouped by domain if Mode 1 is selected
\end{enumerate}

\begin{figure}[h]
\centering
\includegraphics[width=\textwidth]{ui_workflow_step1.png}
\caption{Step 1: URL Entry and Domain Grouping}
\label{fig:ui_workflow_step1}
\end{figure}

\subsection{Step 2: Configure Analysis}

\begin{enumerate}
\item Select DOM analysis levels using the numbered chips (1-10)
\item Choose processing mode:
\begin{itemize}
\item Mode 0: All Together - processes all URLs as a single batch
\item Mode 1: Domain-based - groups URLs by domain for analysis
\end{itemize}
\end{enumerate}

\begin{figure}[h]
\centering
\includegraphics[width=\textwidth]{ui_workflow_step2.png}
\caption{Step 2: Level Selection and Mode Configuration}
\label{fig:ui_workflow_step2}
\end{figure}

\subsection{Step 3: Execute Analysis}

\begin{enumerate}
\item Click the "Run" button to start processing
\item Monitor progress through the loading indicators
\item Processing time varies based on number of URLs and selected levels
\end{enumerate}

\begin{figure}[h]
\centering
\includegraphics[width=\textwidth]{ui_workflow_step3.png}
\caption{Step 3: Analysis Execution with Progress Indicators}
\label{fig:ui_workflow_step3}
\end{figure}

\subsection{Step 4: View Results}

\begin{enumerate}
\item Navigate to the Clusters tab (🎯 Clusters) to view results
\item Results are organized by DOM level
\item Each cluster is color-coded for easy identification
\item Click on sections to expand detailed views
\end{enumerate}

\begin{figure}[h]
\centering
\includegraphics[width=\textwidth]{ui_workflow_step4.png}
\caption{Step 4: Viewing Color-coded Clustering Results}
\label{fig:ui_workflow_step4}
\end{figure}

\section{Advanced Features}
\label{app:advanced_features}

\subsection{Batch Experiments}

For research purposes, use the Experiments tab (🧪 Experiments):

\begin{figure}[h]
\centering
\includegraphics[width=\textwidth]{ui_advanced_experiments.png}
\caption{Advanced Experiments Interface for Batch Processing}
\label{fig:ui_advanced_experiments}
\end{figure}

\subsection{Historical Results}

Access previous analyses through the Results tab (📊 Results):

\begin{figure}[h]
\centering
\includegraphics[width=\textwidth]{ui_advanced_results.png}
\caption{Historical Results Archive with Performance Metrics}
\label{fig:ui_advanced_results}
\end{figure}

\subsection{Custom Configuration}

Fine-tune analysis parameters in the Config tab (⚙️ Config):

\begin{figure}[h]
\centering
\includegraphics[width=\textwidth]{ui_advanced_config.png}
\caption{Custom Configuration for Advanced Users}
\label{fig:ui_advanced_config}
\end{figure}

\chapter{Technical Specifications}
\label{appendix:tech-specs}

\section{System Requirements}
\label{app:system_requirements}

\subsection{Hardware Requirements}

\begin{itemize}
\item \textbf{CPU:} Intel i5-6500 or AMD Ryzen 5 2600 (minimum), Intel i7-8565U or higher (recommended)
\item \textbf{RAM:} 8GB minimum, 16GB recommended for processing large datasets
\item \textbf{Storage:} 2GB available disk space, SSD recommended for optimal performance
\item \textbf{Network:} Stable internet connection for web scraping operations
\end{itemize}

\subsection{Software Requirements}

\begin{itemize}
\item \textbf{Operating System:} Windows 10+, macOS 10.14+, or Linux (Ubuntu 18.04+ recommended)
\item \textbf{Python:} Version 3.8 or higher
\item \textbf{Node.js:} Version 14.0 or higher
\item \textbf{Chrome Browser:} Version 90+ (for ChromeDriver compatibility)
\item \textbf{Git:} For version control and project cloning
\end{itemize}

\section{Installation Guide}
\label{app:installation_guide}

\subsection{Quick Start Installation}

\begin{enumerate}
\item Clone the repository:
\begin{lstlisting}[language=bash]
git clone https://github.com/user/websrc-project.git
cd websrc-project
\end{lstlisting}

\item Install Python dependencies:
\begin{lstlisting}[language=bash]
pip install -r requirements.txt
\end{lstlisting}

\item Install Node.js dependencies:
\begin{lstlisting}[language=bash]
cd frontend
npm install
cd ..
\end{lstlisting}

\item Start the application:
\begin{lstlisting}[language=bash]
python run_project.py
\end{lstlisting}
\end{enumerate}

\section{Configuration Options}
\label{app:configuration}

\subsection{Cache Configuration}

The cache system can be configured through \texttt{backend/lib/config.json}:

\begin{lstlisting}[language=json]
{
  "cache": {
    "enabled": true,
    "default_ttl_hours": 24,
    "max_cache_size_mb": 1024,
    "cleanup_interval_hours": 6
  },
  "clustering": {
    "min_cluster_size_range": [2, 10],
    "cluster_selection_epsilon_range": [0.0, 0.5],
    "outlier_reassignment": true
  }
}
\end{lstlisting}

\subsection{Processing Parameters}

Key processing parameters that can be adjusted:

\begin{table}[h]
\centering
\caption{Processing Configuration Parameters}
\begin{tabular}{|l|l|l|}
\hline
\textbf{Parameter} & \textbf{Default} & \textbf{Description} \\
\hline
\texttt{analysis\_levels} & [1, 2, 3] & DOM depth levels to analyze \\
\texttt{timeout\_seconds} & 30 & Web scraping timeout \\
\texttt{max\_urls\_per\_batch} & 50 & Maximum URLs per processing batch \\
\texttt{enable\_css\_analysis} & true & Include computed CSS styles \\
\texttt{dbcv\_optimization} & true & Enable DBCV-based parameter optimization \\
\hline
\end{tabular}
\end{table}

\section{Performance Tuning Guide}
\label{app:performance_tuning}

\subsection{Memory Optimization}

\begin{itemize}
\item Adjust \texttt{max\_urls\_per\_batch} based on available RAM
\item Enable automatic garbage collection for large datasets
\item Use streaming processing for very large URL lists
\item Monitor memory usage during processing
\end{itemize}

\subsection{CPU Optimization}

\begin{itemize}
\item Configure parallel processing based on CPU cores
\item Optimize clustering parameter ranges for faster convergence
\item Use browser pooling for multiple concurrent scraping operations
\item Enable CPU profiling for performance bottleneck identification
\end{itemize}

\section{Troubleshooting Guide}
\label{app:troubleshooting}

\subsection{Common Issues}

\subsubsection{ChromeDriver Compatibility}

\textbf{Issue:} ChromeDriver version mismatch with installed Chrome browser.

\textbf{Solution:}
\begin{lstlisting}[language=bash]
# Update Chrome browser to latest version
# Install compatible ChromeDriver
pip install --upgrade selenium webdriver-manager
\end{lstlisting}

\subsubsection{Memory Errors}

\textbf{Issue:} Out of memory errors during large dataset processing.

\textbf{Solution:}
\begin{itemize}
\item Reduce \texttt{max\_urls\_per\_batch} parameter
\item Increase system virtual memory
\item Process datasets in smaller chunks
\item Enable streaming processing mode
\end{itemize}

\subsubsection{Cache Corruption}

\textbf{Issue:} Corrupted cache files causing processing errors.

\textbf{Solution:}
\begin{lstlisting}[language=bash]
# Clear cache directory
rm -rf backend/api/cache/*
# Restart application
python run_project.py
\end{lstlisting} 